\documentclass[11pt,a4paper]{article}
\usepackage[margin=1in]{geometry}
\usepackage{lmodern}
\usepackage[T1]{fontenc}
\usepackage{microtype}
\usepackage{hyperref}
\usepackage{enumitem}
\usepackage{graphicx}
\usepackage{xcolor}

% Give LaTeX more room to stretch/shrink lines before
% resorting to an overfull box
\setlength{\emergencystretch}{3em}

% Variable: Lab Instructor
\makeatletter \newcommand{\BvrSetLabInstructor}[1]{%
  \gdef\Bvr@LabInstructor{#1}}
% default
\newcommand{\Bvr@LabInstructor}{Dr. Paramveer Sidhu}
% public accessor
\newcommand{\BvrLabInstructorValue}{\Bvr@LabInstructor}
\makeatother

% Add subtitle
\usepackage{titling}
% -- Set title bold
\pretitle{\begin{center}\bfseries\fontsize{18bp}{18bp}\selectfont}
\makeatletter
\newcommand{\BvrSubtitle}[1]{%
  \posttitle{%
    \par\end{center}
    \begin{center}\large#1\end{center}
    \vskip 0.5em}%
}
\makeatother

% -----------------------------------------------------

\title{Driver Drowsiness Detection and Alert System}

\BvrSetLabInstructor{Dr.Paramveer Sidhu}

\BvrSubtitle{%
  Project Proposal for UCS503  \\
  Submitted to: \BvrLabInstructorValue \\ \bfseries
  Thapar Institute of Engineering and Technology%
}

\author{
   Arshita Gupta, 3C53
  \thanks{Roll No: \texttt{1024030725}, 
    Email: \texttt{agupta24\_be24\@thapar.edu}}
  \and
   Tarushi Goyal, 3C53%
  \thanks{Roll No: \texttt{1024030730}, %
    Email: \texttt{tgoyal2\_be24\@thapar.edu}}
     \and
   Chahana Goyal, 3C53%
  \thanks{Roll No: \texttt{1024030746}, %
    Email: \texttt{cgoyal\_be24\@thapar.edu}}
}

\date{\today}

\begin{document}
\maketitle

\section{Introduction}
This project aims to improve road safety by reducing
accidents caused by driver fatigue and micro-sleep
episodes. While the underlying concern (fatigue-related
crashes) is a broader public-health and transportation
issue, the immediate engineering objective is to
translate \emph{real-time drowsiness detection} into a
measurable, deployable software solution running on
commodity hardware (webcam/smartphone camera).

\section{Objectives}

The primary objective of this project is to develop a lightweight,
real-time camera-based system for detecting driver drowsiness and
providing timely alerts.

The specific objectives are:
\begin{itemize}[leftmargin=0.7cm]
    \item To detect the driver's face and facial landmarks in real time
    using a webcam or dashboard-mounted camera.
    
    \item To detect prolonged eye closure and yawning using Eye Aspect
    Ratio (EAR) and Mouth Aspect Ratio (MAR).
    
    \item To reduce false alarms using temporal smoothing and
    per-user calibration.
    
    \item To provide an audio/visual alert when sustained signs of
    drowsiness are detected.
    
    \item To evaluate the system using alert latency, detection
    accuracy, false-alarm rate, and real-time performance.
\end{itemize}

\section{Problem Statement}
Drowsy driving is a significant contributor to road
accidents, particularly on highways and during long or
night-time drives. Common pain points include:
\begin{itemize}[leftmargin=0.7cm]
    \item \textbf{Late detection:} drivers and passengers often fail to notice early signs of fatigue (prolonged blinks, yawning, head nodding) until it is too late.
    \item \textbf{Cost barrier:} existing dedicated drowsiness-detection hardware (e.g., specialized in-cabin sensors) is expensive and not standard in most vehicles, especially in tier-2/3 markets.
    \item \textbf{False alarms:} naive systems using rigid thresholds trigger frequent false positives (e.g., due to glasses, lighting, or eye shape variation), reducing driver trust and adoption.
    \item \textbf{Lack of accessible logging:} there is no easy way for fleet operators or individual users to review drowsiness incidents after a trip.
\end{itemize}

\textbf{Why this matters now (contextual relevance):}
With cameras already available in smartphones and
low-cost dashboard-mounted devices, a software-only or
software-plus-commodity-camera solution can meaningfully
reduce fatigue-related risk without expensive
specialized hardware.

\section{Proposed Solution}
\subsection{Overview}
Build a \textbf{camera-based} ``Drowsiness-to-Alert''
system with the following components:
\begin{itemize}[leftmargin=0.7cm]
    \item \textbf{Capture module:} continuous video capture from a webcam or dashboard-mounted camera.
    \item \textbf{Face and landmark detection:} locate the driver's face and extract eye/mouth landmarks per frame.
    \item \textbf{Drowsiness scoring:} compute indicators such as Eye Aspect Ratio (EAR) and Mouth Aspect Ratio (MAR) to flag prolonged eye closure and yawning.
    \item \textbf{Alert mechanism:} trigger an audio/visual alert when a temporally-smoothed drowsiness score crosses a threshold.
    \item \textbf{Event logging dashboard:} record timestamped drowsiness events (without storing raw video) for later review.
\end{itemize}

\subsection{Core workflow}
\begin{enumerate}[leftmargin=0.7cm]
    \item Camera captures a frame; the system detects the driver's face and facial landmarks.
    \item EAR and MAR are computed per frame; a rolling window smooths noisy per-frame values.
    \item If eye closure or yawning persists beyond a calibrated duration, a drowsiness state is flagged.
    \item An alert (sound/vibration/visual) is triggered, and the event is logged with a timestamp and severity level.
\end{enumerate}

\subsection{Operational constraints for an educational,
  real-time setting}
\begin{itemize}[leftmargin=0.7cm]
    \item Keep the pipeline lightweight enough to run in real time (target $\geq$ 15 FPS) on low-to-mid range hardware (e.g., laptop webcam, Raspberry Pi, or a mid-range smartphone).
    \item Avoid dependence on training novel deep-learning models from scratch; use pre-trained facial landmark detectors and well-established heuristics.
    \item Design for deployability as a standalone desktop/mobile app with an optional lightweight backend for logging and fleet-level aggregation.
\end{itemize}

\section{Solution Approach}
This is an engineering course project, so we focus on
integrating mature computer-vision components into a
reliable, real-time, measurable pipeline rather than
researching new detection algorithms.

\subsection{Drowsiness scoring}
Detection will be heuristic-based, built on established,
well-documented techniques:
\begin{itemize}[leftmargin=0.7cm]
    \item Use a pre-trained facial landmark model (e.g., dlib's 68-point model or MediaPipe Face Mesh) rather than training a new detector.
    \item Compute EAR (eye aspect ratio) and MAR (mouth aspect ratio) using standard published formulas, with calibrated per-user thresholds.
    \item Apply temporal smoothing (e.g., a sliding window / consecutive-frame count) to reduce false positives from single-frame noise, and surface borderline cases as ``low-confidence'' alerts rather than silently suppressing them.
\end{itemize}

\subsection{System architecture}
A layered architecture with an edge-first design:
\begin{itemize}[leftmargin=0.7cm]
    \item \textbf{Edge/client layer:} camera capture, face/landmark detection, scoring, and alerting run locally to minimize latency and avoid transmitting raw video.
    \item \textbf{Backend API (optional):} receives only lightweight event metadata (timestamp, severity, session ID) for logging and analytics; no raw frames are uploaded by default.
    \item \textbf{Database:} stores users, sessions, drowsiness events, and calibration profiles.
\end{itemize}

\subsection{CI/CD and fast delivery enabler}
CI/CD will be used to:
\begin{itemize}[leftmargin=0.7cm]
    \item Automatically run unit tests on the scoring logic (EAR/MAR computation, thresholding, smoothing) on every change.
    \item Run pipeline tests against a fixed set of labeled sample videos to catch regressions in detection behavior.
    \item Produce repeatable deployment artifacts for the desktop/mobile client and backend to staging.
\end{itemize}

\section{Evaluation Criterion}
We define success using measurable metrics tied to
detection accuracy and real-time responsiveness.

\subsection{Primary evaluation metric}
\textbf{Alert latency:} time between the onset of a
drowsiness event (e.g., sustained eye closure) in the
input video and the system raising an alert.
\begin{itemize}[leftmargin=0.7cm]
    \item Target: median alert latency $\le$ 1 second on labeled test clips.
    \item Attribution: measured by comparing ground-truth event onset timestamps (from labeled/annotated test videos) against logged alert timestamps.
\end{itemize}

\subsection{Secondary evaluation metrics}
\begin{itemize}[leftmargin=0.7cm]
    \item \textbf{Detection accuracy:} sensitivity (recall) and specificity of drowsiness classification against a labeled validation set (e.g., a subset of an existing public drowsiness dataset or self-recorded, consented clips).
    \item \textbf{False alarm rate:} number of false-positive alerts per hour of normal (alert) driving footage.
    \item \textbf{Real-time performance:} sustained frames-per-second on target hardware, and CPU/memory usage.
    \item \textbf{User-perceived reliability:} a simple 1--5 feedback question on alert timeliness and annoyance level from pilot users.
\end{itemize}

\subsection{Pilot validation plan}
\begin{itemize}[leftmargin=0.7cm]
    \item Collect or use an existing small set of labeled driving/alertness video clips (with consent, or a public dataset) covering both alert and simulated-drowsy states.
    \item Run 5--10 pilot users through short supervised sessions (e.g., a driving simulator or seated webcam test) to measure real-world alert latency and false-alarm rate.
\end{itemize}

\section{Timeline}

\begin{center}
\renewcommand{\arraystretch}{1.3}
\begin{tabular}{|c|p{6cm}|p{5cm}|}
\hline
\textbf{Weeks} & \textbf{Task} & \textbf{Deliverable} \\
\hline

1--2 &
Literature review, requirement analysis, and project setup &
Requirements specification and technology stack \\
\hline

3--4 &
Camera capture and facial landmark detection &
Working face and landmark detection module \\
\hline

5--6 &
EAR-based eye-closure detection and threshold calibration &
Eye-closure detection prototype \\
\hline

7--8 &
MAR-based yawning detection and combined drowsiness scoring &
Drowsiness scoring module \\
\hline

9--10 &
Alert mechanism and event logging integration &
Working end-to-end detection and alert system \\
\hline

11 &
Testing, CI/CD integration, and performance evaluation &
Test results and performance report \\
\hline

12 &
Buffer week, refinement, documentation, and final testing &
Final prototype and project documentation \\
\hline

\end{tabular}
\end{center}

\section{Scalability}
The proposition should scale in both theoretical and
practical senses.

\begin{enumerate}[leftmargin=0.7cm]
    \item \textbf{Scalable (theoretically):} The system should support growth from single-driver monitoring to fleet-level analytics by:
    \begin{itemize}[leftmargin=0.7cm]
        \item keeping heavy processing on the edge device so backend load only scales with event metadata, not video volume,
        \item using stateless, indexed APIs for querying historical events across many drivers/vehicles,
        \item allowing per-user calibration profiles to be stored and retrieved independently.
    \end{itemize}

    \item \textbf{Operationally deployable (practically):} The system should run on common hardware and hosting:
    \begin{itemize}[leftmargin=0.7cm]
        \item lightweight client for laptops, Raspberry Pi-class devices, or smartphones,
        \item managed database and containerized backend for the optional logging service,
        \item CI/CD pipeline for staging/production releases of both client and backend.
    \end{itemize}
\end{enumerate}

\section{Technology Stack}
A mature set of ``engines'' (tooling/libraries) is
available to implement this without inventing new
computer-vision research:
\begin{itemize}[leftmargin=0.7cm]
    \item Computer-vision library for frame capture and preprocessing (OpenCV).
    \item Pre-trained facial landmark detector (dlib 68-point model or MediaPipe Face Mesh).
    \item Audio/notification library for triggering alerts on the client device.
    \item Lightweight web/mobile framework for the client UI and optional dashboard.
    \item Managed relational or document database for event logging.
    \item Test framework for unit/integration tests on scoring logic.
    \item CI runner and staging deployment target.
\end{itemize}
We will use these mature components to focus on
pipeline integration, calibration, correctness, and
measurement rather than building detection models from
scratch.

\section{Project scope and deliverables}
\subsection{Initial deliverable}
\begin{itemize}[leftmargin=0.7cm]
    \item Webcam capture + facial landmark detection integration
    \item EAR-based eye-closure detection with basic thresholding and temporal smoothing
    \item Local audio/visual alert on sustained eye closure
    \item Basic event logging with timestamps for latency measurement
    \item CI: build + unit tests for scoring logic + linting
    \item CD to staging for the client application
\end{itemize}

\subsection{Subsequent deliverables}
\begin{itemize}[leftmargin=0.7cm]
    \item MAR-based yawning detection and combined drowsiness scoring
    \item Per-user calibration flow to reduce false alarms
    \item Optional backend for event logging and a simple analytics dashboard
    \item Low-light/IR-camera considerations for night-time driving scenarios
\end{itemize}

\section{Risks and mitigations}
\begin{itemize}[leftmargin=0.7cm]
    \item \textbf{Privacy of video data:} process video entirely on-device where possible; never store or transmit raw frames by default, only derived event metadata.
    \item \textbf{False positives/negatives:} mitigate via per-user calibration, temporal smoothing, and clearly logging low-confidence events instead of silently discarding or over-alerting.
    \item \textbf{Hardware/lighting variability:} test across multiple lighting conditions and camera qualities early; document known limitations (e.g., sunglasses, extreme low light).
    \item \textbf{Driver distraction from alerts:} keep alert design (sound/vibration pattern) simple and non-startling, validated informally with pilot users.
    \item \textbf{Operational issues in pilot:} use staging and recorded test videos early; ensure CI/CD produces repeatable, testable builds before live testing.
\end{itemize}

\section{Summary}
This project proposes a deployable, camera-based
drowsiness prediction and alerting system to reduce
fatigue-related driving risk. It meets the course
criteria by emphasizing time-to-value through rapid
iterations built on pre-trained, mature computer-vision
components, implementing CI/CD to support frequent safe
releases, and using measurable evaluation criteria
(especially alert latency and detection accuracy). It
remains focused on engineering integration, real-time
reliability, and scalability readiness rather than
novel algorithmic research.

\end{document}